% !TeX encoding = UTF-8
% !TeX program = pdflatex
% ---------------------------------------------------------------------------
% IEEE TGRS (Transactions on Geoscience and Remote Sensing) Journal Template
%
% Usage:
%   1. Requires IEEEtran.cls and IEEEtran.bst (included in this directory)
%   2. Compile: pdflatex template.tex → bibtex template → pdflatex template.tex → pdflatex template.tex
%
% Journal: IEEE Transactions on Geoscience and Remote Sensing (TGRS)
% Publisher: IEEE Geoscience and Remote Sensing Society (GRSS)
% ISSN: 0196-2892
% Page limit: typically 12-18 pages (no strict limit for initial submission)
% ---------------------------------------------------------------------------

\documentclass[journal,onecolumn,draftcls]{IEEEtran}
\IEEEoverridecommandlockouts

% --- Switch to final version for camera-ready ---
% \documentclass[journal]{IEEEtran}
% \IEEEoverridecommandlockouts

\usepackage{cite}
\usepackage{amsmath,amssymb,amsfonts}
\usepackage{algorithmic}
\usepackage{algorithm}
\usepackage{graphicx}
\usepackage{textcomp}
\usepackage{xcolor}
\usepackage{booktabs}
\usepackage{multirow}
\usepackage{arydshln}  % dashed lines in tables
\usepackage{hyperref}
\usepackage{gensymb}   % \degree, \celsius, etc. (important for remote sensing)
\usepackage{siunitx}   % SI units: \SI{10}{\meter}
\usepackage{subcaption}

\hypersetup{
  colorlinks=true,
  linkcolor=blue,
  citecolor=blue,
  urlcolor=blue,
}

\def\BibTeX{{\rm B\kern-.05em{\sc i\kern-.025em b}\kern-.08em
    T\kern-.1667em\lower.7ex\hbox{E}\kern-.125emX}}

% --- TGRS Journal Metadata ---
% \markboth{IEEE Transactions on Geoscience and Remote Sensing, Vol. XX, No. X, Month 2026}%
% {Author \MakeLowercase{\textit{et al.}}: Short Title Here}

\begin{document}

% ===========================================================================
% TITLE
% ===========================================================================
\title{Your Paper Title: A Comprehensive Study of [Method] for [Remote Sensing Task]}

\author{
  \IEEEauthorblockN{Author One\IEEEauthorrefmark{1},
                    Author Two\IEEEauthorrefmark{2},
                    Author Three\IEEEauthorrefmark{1}}
  \IEEEauthorblockA{
    \IEEEauthorrefmark{1}Department of XXX, University of YYY, City, Country \\
    \IEEEauthorrefmark{2}Institute of ZZZ, Academy of Sciences, City, Country \\
    E-mail: \{email1, email2\}@university.edu
  }
}

% TGRS typically allows long titles. Use descriptive titles that include:
% - The remote sensing task (e.g., classification, detection, change detection)
% - The sensor/modality (e.g., SAR, multispectral, hyperspectral, LiDAR)
% - The core technique (e.g., transformer, diffusion, contrastive learning)

\maketitle

% ===========================================================================
% ABSTRACT
% ===========================================================================
\begin{abstract}

% TGRS abstracts: typically 150-300 words in a single paragraph.
% Follow modified 2-6-2 structure adapted for journal length:
%   2-3 sentences: background + remote sensing context + specific gap
%   5-8 sentences: method (First/Second/Finally)
%   2-3 sentences: results + broader impact for geoscience applications

Your abstract here. For TGRS, emphasize:
\begin{enumerate}
  \item The specific remote sensing application and its practical importance
  \item The sensor modality (optical/SAR/multispectral/hyperspectral/LiDAR/GNSS)
  \item The gap in existing geoscience methods
  \item Your method's key innovations
  \item Quantitative results on standard remote sensing benchmarks
\end{enumerate}

\textbf{Example:} Remote sensing image classification is essential for land
cover mapping, urban planning, and environmental monitoring. However, existing
deep learning methods struggle with [specific challenge, e.g., "the inherent
class imbalance between dominant and rare land cover types in large-scale
satellite imagery"]. To address this, we propose [Method], which [core mechanism].
First, [component 1 + why it works]. Second, [component 2 + why it works].
Finally, [component 3 + why it works]. Experiments on [N] benchmark datasets
([Dataset1, Dataset2, Dataset3]) demonstrate that [Method] improves [metric] by
[X]\% over [strongest baseline], with particularly strong gains on
underrepresented classes (+[Y]\%).
\end{abstract}

% ===========================================================================
% KEYWORDS — TGRS uses IEEE keywords
% ===========================================================================
\begin{IEEEkeywords}
  Remote sensing, [your task, e.g., land cover classification],
  [your modality, e.g., multispectral imagery],
  [your technique, e.g., deep learning, attention mechanism],
  [secondary technique, e.g., self-supervised learning].
\end{IEEEkeywords}

% ===========================================================================
% 1. INTRODUCTION
% ===========================================================================
\section{Introduction}
\label{sec:introduction}

% TGRS Introduction: typically longer than conference papers (~2-3 pages in
% the submitted manuscript). Structure in 5 paragraphs expanded for journal depth:

% Para 1: Remote sensing context — why this task matters for geoscience.
%   - Cite specific applications: land cover mapping, disaster monitoring,
%     agricultural assessment, urban growth, climate change analysis, etc.
%   - Connect to Earth observation programs (Landsat, Sentinel, MODIS, etc.)

% Para 2: Technical progress in the task.
%   - Review 2-3 major method families with specific papers.
%   - Track the evolution: traditional ML → CNN → ViT/Transformer → foundation models.
%   - Note the unique challenges of remote sensing data: large spatial extent,
%     multi-source heterogeneity, spectral variability, temporal dynamics.

% Para 3: Remaining gap — what current geoscience methods fail to address.
%   - Must contain a "because" clause specific to remote sensing.
%   - Example: "Existing methods assume X, but in remote sensing Y occurs because Z."

% Para 4: Proposed method teaser — 3-4 sentences, core intuition.
%   - Explain WHY your method is suited for remote sensing data specifically.

% Para 5: Main contributions (4-5 bullets for journal papers):
%   - Each bullet should map to a specific Table/Figure/Section.
%   - TGRS reviewers expect at least one contribution to be validated on
%     multiple remote sensing datasets and/or multiple sensor modalities.

% Example contribution paragraph structure:
The main contributions of this paper are as follows:

\begin{enumerate}
  \item We identify [specific limitation] in existing [method family] when
    applied to [remote sensing scenario], and we validate this limitation
    quantitatively (Section~\ref{sec:problem-analysis}).
  \item We propose [Method Name], a [brief description] that [core mechanism]
    tailored for [remote sensing data characteristic]
    (Section~\ref{sec:method}).
  \item We design [specific component], which [what it does] to address
    [remote sensing challenge] (Section~\ref{sec:method-component}).
  \item Through extensive experiments on [N] datasets spanning [sensor types
    / resolutions / geographic regions], we demonstrate that [Method] achieves
    [key result], outperforming [strongest baseline] by [X]\%
    (Section~\ref{sec:experiments}).
  \item We release our code and [trained models / dataset splits] to facilitate
    reproducible research in remote sensing (URL: to be added).
\end{enumerate}

% ===========================================================================
% 2. RELATED WORK
% ===========================================================================
\section{Related Work}
\label{sec:related}

% TGRS Related Work is typically more comprehensive than conference papers.
% Organize by 4-5 themes, each with detailed discussion.

\subsection{Remote Sensing [Task]}
\label{sec:related-task}

% Review task-specific methods: traditional approaches → deep learning evolution.
% Cite TGRS, GRSL, JSTARS, RSE, ISPRS journal papers to show domain expertise.

\subsection{[Method Family A]}
\label{sec:related-familyA}

% Discuss the method family your approach builds on or competes with.

\subsection{[Method Family B]}
\label{sec:related-familyB}

% Second relevant method family.

\subsection{Earth Observation Foundation Models}
\label{sec:related-foundation}

% If relevant, discuss recent geospatial foundation models:
% - SatMAE, Scale-MAE, GeoViT, Prithvi, SatCLIP, GFM, etc.
% - Compare to general-purpose models (CLIP, DINOv2) adapted for remote sensing.

\subsection{Summary and Positioning}
\label{sec:related-summary}

% 2-3 sentence summary: what distinguishes your work from all reviewed prior work.
% Be specific about the gap your method fills.

% ===========================================================================
% 3. PROBLEM FORMULATION (optional but common in TGRS)
% ===========================================================================
\section{Problem Formulation}
\label{sec:problem}

% TGRS papers often include a formal problem definition section.
% Define notation, input/output spaces, assumptions, and evaluation criteria.

\subsection{Notation and Definitions}
\label{sec:problem-notation}

% Define key symbols: input data X (e.g., multispectral image patches),
% labels Y, sensor characteristics, spatial resolution, spectral bands, etc.

\subsection{Task Definition}
\label{sec:problem-task}

% Formal definition of the remote sensing task your paper addresses.

% ===========================================================================
% 4. PROPOSED METHOD
% ===========================================================================
\section{Proposed Method}
\label{sec:method}

% TGRS method sections should be detailed enough for reproducible research.
% Aim for ~3-5 pages with clear sub-sections.

\subsection{Overview}
\label{sec:method-overview}

% Architecture diagram: Fig. 1. Describe the complete pipeline.
% For remote sensing: emphasize how your method handles unique RS data
% characteristics (large coverage, multi-scale, multi-source, temporal).

\begin{figure}[t]
  \centering
  \fbox{\parbox{0.85\linewidth}{\centering
  \vspace{3cm}
  {\large [Fig. 1: Overall architecture of the proposed method.]}\\
  \vspace{0.3cm}
  {\footnotesize The pipeline shows: (a) [Input preprocessing],
   (b) [Feature extraction backbone],
   (c) [Core innovation module],
   (d) [Output/task head].}
  \vspace{3cm}}}
  \caption{Overview of the proposed [Method Name] architecture.
           [One-sentence description of the data flow].}
  \label{fig:architecture}
\end{figure}

\subsection{[Core Innovation Module A]}
\label{sec:method-moduleA}

% For each module: Input → Process → Output → Rationale.
% Include equations with ALL symbols defined.
% Discuss WHY this module is specifically suited for remote sensing data.

\subsection{[Core Innovation Module B]}
\label{sec:method-moduleB}

% Second key innovation.

\subsection{Loss Function and Training Strategy}
\label{sec:method-loss}

% Define all loss terms. If multi-task, explain task balancing.
% Remote sensing specific: class imbalance handling, spatial consistency
% regularization, spectral-photometric invariance, etc.

\begin{equation}
  \mathcal{L}_{\text{total}} =
    \mathcal{L}_{\text{task}} +
    \lambda_{1} \mathcal{L}_{\text{aux1}} +
    \lambda_{2} \mathcal{L}_{\text{aux2}}
  \label{eq:total_loss}
\end{equation}

\subsection{Implementation Details}
\label{sec:method-implementation}

% Comprehensive for journal reproducibility:
% - Framework (PyTorch version, CUDA version)
% - Optimizer type + hyperparameters
% - Learning rate schedule + warmup
% - Batch size + gradient accumulation
% - Hardware (GPU model, count, memory)
% - Training duration (wall-clock time, epochs)
% - Data preprocessing pipeline
% - Data augmentation (geometric + spectral)
% - Pretrained weights used (if any)

% ===========================================================================
% 5. EXPERIMENTS
% ===========================================================================
\section{Experiments}
\label{sec:experiments}

% TGRS experiments: typically very comprehensive (4-6 pages).
% Expected components: datasets, implementation, main results, ablation,
% parameter sensitivity, computational efficiency, qualitative visualization,
% failure case analysis.

% ---------------------------------------------------------------------------
\subsection{Datasets}
\label{sec:exp-datasets}

% For each dataset, describe:
% - Sensor/platform (e.g., Sentinel-2, Landsat-8, Gaofen-2, WorldView-3)
% - Spatial resolution and coverage
% - Spectral bands and wavelengths
% - Temporal coverage (if multi-temporal)
% - Geographic region
% - Training/validation/test split
% - Number of samples per class
% - Any preprocessing (cloud removal, atmospheric correction, etc.)

\begin{table}[t]
  \centering
  \caption{Summary of datasets used in experiments. GSD = Ground Sampling Distance.}
  \label{tab:datasets}
  \begin{tabular}{lccccr}
    \toprule
    \textbf{Dataset} & \textbf{Sensor} & \textbf{GSD (m)} & \textbf{Bands} &
    \textbf{Classes} & \textbf{Images} \\
    \midrule
    Dataset 1 & Sensor A & 10 & 13 (MS) & 15 & 5,000 \\
    Dataset 2 & Sensor B & 0.5 & 3 (RGB) & 8  & 2,000 \\
    Dataset 3 & Sensor C & 30 & 6 (MS) & 10 & 1,500 \\
    \bottomrule
  \end{tabular}
\end{table}

% ---------------------------------------------------------------------------
\subsection{Evaluation Metrics}
\label{sec:exp-metrics}

% TGRS common metrics:
% Classification: OA (Overall Accuracy), AA (Average Accuracy), Kappa (κ),
%                 per-class F1, confusion matrix
% Segmentation: mIoU, FWIoU, per-class IoU
% Detection: mAP, AP@[.5:.95]
% Change detection: OA, F1, Kappa
% Regression: RMSE, MAE, R²

% Define each metric and state direction (higher/lower is better).

% ---------------------------------------------------------------------------
\subsection{Comparison with State-of-the-Art}
\label{sec:exp-main}

% Main results table: bold = best, underline = second best.
% Include: method name, year, OA, AA, Kappa, per-class metrics, params, FLOPs.
% Use \pm for standard deviations over multiple runs.

\begin{table}[t]
  \centering
  \caption{Comparison with state-of-the-art methods on [Dataset].
           Best results in \textbf{bold}, second-best \underline{underlined}.
           All results are mean $\pm$ std over 5 runs.}
  \label{tab:main_results}
  \begin{tabular}{lccccc}
    \toprule
    \textbf{Method} & \textbf{OA (\%)} & \textbf{AA (\%)} &
    \textbf{$\kappa \times 100$} & \textbf{Params (M)} & \textbf{FLOPs (G)} \\
    \midrule
    SVM \cite{ref1}           & 82.3 $\pm$ 0.5 & 79.1 $\pm$ 0.8 &
    0.801 & -- & -- \\
    CNN-2D \cite{ref2}        & 90.5 $\pm$ 0.3 & 88.7 $\pm$ 0.5 &
    0.895 & 5.2 & 8.3 \\
    ResNet-50 \cite{ref3}     & 92.1 $\pm$ 0.2 & 90.3 $\pm$ 0.4 &
    0.913 & 23.5 & 4.1 \\
    ViT-B \cite{ref4}         & 93.2 $\pm$ 0.3 & 91.5 $\pm$ 0.4 &
    0.925 & 86.0 & 17.6 \\
    \textbf{Ours}             & \textbf{94.8 $\pm$ 0.2} &
    \textbf{93.2 $\pm$ 0.3} & \textbf{0.942} & 28.1 & 5.3 \\
    \bottomrule
  \end{tabular}
\end{table}

% ---------------------------------------------------------------------------
\subsection{Ablation Study}
\label{sec:exp-ablation}

% Comprehensive ablation: each claimed component tested individually.
% Use waterfall-style table (full model → remove component A → ... → baseline).

% ---------------------------------------------------------------------------
\subsection{Parameter Sensitivity Analysis}
\label{sec:exp-sensitivity}

% TGRS reviewers expect sensitivity analysis for key hyperparameters.
% Show how performance varies with λ, temperature, window size, etc.
% Use line plots (Fig. X) showing stability range.

% ---------------------------------------------------------------------------
\subsection{Computational Efficiency}
\label{sec:exp-efficiency}

% For TGRS: efficiency matters for large-scale remote sensing deployment.
% Report: training time, inference time (per image / per km²), GPU memory,
% model size, throughput.

% ---------------------------------------------------------------------------
\subsection{Qualitative Analysis and Visualization}
\label{sec:exp-qualitative}

% Show classification maps, attention heatmaps, feature visualizations.
% Compare predicted maps against ground truth.
% Show BOTH successes and failure cases — TGRS reviewers value honesty.

% ===========================================================================
% 6. DISCUSSION
% ===========================================================================
\section{Discussion}
\label{sec:discussion}

% TGRS journal papers often have a dedicated Discussion section.
% Topics to cover:
\subsection{Key Findings and Implications}
\label{sec:discuss-findings}

\subsection{Why the Method Works for Remote Sensing Data}
\label{sec:discuss-why}

\subsection{Limitations and Failure Cases}
\label{sec:discuss-limitations}

% Be specific about:
% - Geographic transferability (does it work in different regions?)
% - Sensor transferability (does it work with different sensors?)
% - Temporal robustness (does it work across seasons?)
% - Computational scalability (can it handle country-scale mapping?)

\subsection{Future Work}
\label{sec:discuss-future}

% Concrete directions, not generic statements.
% Example: "Extending to [other sensor modality, e.g., SAR] by [specific approach]."
% Example: "Integrating physical models ([specific model, e.g., PROSAIL]) as a prior."

% ===========================================================================
% 7. CONCLUSION
% ===========================================================================
\section{Conclusion}
\label{sec:conclusion}

% Restate: problem, method, key findings. NO new claims or citations.
% 2-3 paragraphs for journal length.

% ===========================================================================
% ACKNOWLEDGMENTS
% ===========================================================================
\section*{Acknowledgments}

% Acknowledge funding, data providers (ESA/NASA/USGS for satellite data),
% computing resources.

% ===========================================================================
% APPENDICES (if needed, placed after references)
% ===========================================================================
% \appendices
% \section{Derivation of [Equation]}
% \section{Additional Experimental Results}

% ===========================================================================
% REFERENCES
% ===========================================================================
{
  % TGRS uses numeric IEEE citation style: [1], [2], ...
  % Include ~40-80 references for a journal submission.
  % Balance: classic RS papers, recent deep learning papers, domain journals.

  \bibliographystyle{IEEEtran}
  \bibliography{references}
}

\end{document}
